% !TEX root = main.tex
\documentclass[14pt,a4paper]{extarticle}

\usepackage[utf8]{inputenc}
\usepackage[T2A]{fontenc}
\usepackage[ukrainian]{babel}

\usepackage[
    left=25mm,
    right=10mm,
    top=20mm,
    bottom=20mm
]{geometry}

% Міжрядковий інтервал
\usepackage{setspace}
\onehalfspacing

% Відступ абзацу
\usepackage{indentfirst}
\setlength{\parindent}{1.25cm}

% Математика
\usepackage{amsmath,amssymb,amsthm}

% Графіка
\usepackage{graphicx}
\graphicspath{{images/}}

% Таблиці
\usepackage{booktabs}
\usepackage{longtable}

% Посилання та гіперпосилання
\usepackage{hyperref}
\hypersetup{
    colorlinks=true,
    linkcolor=black,
    filecolor=black,
    urlcolor=blue,
    citecolor=black
}

% Цитування
\usepackage{csquotes}

% Заголовки секцій
\usepackage{titlesec}
\titleformat{\section}
    {\normalfont\bfseries\centering}{\thesection}{1em}{}
\titleformat{\subsection}
    {\normalfont\bfseries}{\thesubsection}{1em}{}

% Нумерація сторінок
\usepackage{fancyhdr}
\pagestyle{fancy}
\fancyhf{}
\fancyhead{}
\fancyhead[R]{\thepage}
\renewcommand{\headrulewidth}{0pt}

% Лістинги коду (якщо потрібно)
\usepackage{listings}
\usepackage{xcolor}

\lstset{
    basicstyle=\ttfamily\small,
    keywordstyle=\color{blue},
    commentstyle=\color{green!50!black},
    stringstyle=\color{red},
    numbers=left,
    numberstyle=\tiny,
    breaklines=true,
    frame=single
}

% ===== ПОЧАТОК ДОКУМЕНТА =====
\begin{document}

% Титульна сторінка
\begin{titlepage}
	\centering
	\textbf{МІНІСТЕРСТВО ОСВІТИ І НАУКИ УКРАЇНИ}\\
	\textbf{ЛЬВІВСЬКИЙ НАЦІОНАЛЬНИЙ УНІВЕРСИТЕТ ІМЕНІ ІВАНА ФРАНКА}\\[0.5cm]
	\begin{tabular}{c}
		{Факультет прикладної математики та інформатики} \\
		{Кафедра математичного моделювання соціально-економічних процесів}
	\end{tabular}

	\vspace{3cm}

	{\Large\textbf{КУРСОВА РОБОТА}}\\[0.5cm]
	{\large\textbf{Дослідження та практична реалізація алгоритмів захисту від 2D-атак в задачах розпізнавання облич}}\\

	\vfill

	\begin{flushright}
		\begin{tabular}{ll}
			Виконав: студент групи \underline{ПМАМ-12}                                                                                                                      \\[0.5ex]
			спеціальності:
			\begin{tabular}[t]{c}
				\underline{F4 - Системний аналіз} \\[-1.5ex]
				\scriptsize (шифр і назва спеціальності)
			\end{tabular}                                                                                                                     \\
			\begin{tabular}[b]{@{}c@{}c@{}c@{}}
				{\makebox[2.3cm]{\raisebox{-0.57ex}{\phantom{Студент}}}} & \underline{\makebox[2cm]{\strut}} & \underline{\makebox[4.3cm]{\raisebox{-0.57ex}{Корицький А. О.}}} \\[-1.1ex]
				\scriptsize\phantom{text}                                & \scriptsize(підпис)               & \scriptsize(прізвище та ініціали)
			\end{tabular} \\
			\begin{tabular}[b]{@{}c@{}c@{}c@{}}
				{\makebox[2.3cm]{\raisebox{-0.8ex}{Керівник}}} & \underline{\makebox[2cm]{\strut}} & \underline{\makebox[4.3cm]{\raisebox{-0.57ex}{Добуляк Л. П.}}} \\[-1.1ex]
				                                               & \scriptsize(підпис)               & \scriptsize(прізвище та ініціали)
			\end{tabular}
		\end{tabular}
	\end{flushright}

	\vfill

	Львів -- 2026
\end{titlepage}
\setcounter{page}{2}

% Зміст
\pagestyle{empty}
\tableofcontents
\newpage
\pagestyle{fancy}

%? Розділ 1: Вступ і огляд проблематики
\section{Вступ і огляд проблематики}
\subsection{Біометрична верифікація та системи розпізнавання облич}
Біометрична верифікація — це процес автоматичного підтвердження особистості людини на основі її унікальних фізіологічних або поведінкових характеристик.
На відміну від традиційних методів автентифікації, таких як паролі або PIN-коди, біометричні системи спираються на властивості самої людини, що унеможливлює
їх забуття чи передачу третім особам \cite{Tome2015}.
Серед усіх біометричних модальностей — відбитків пальців, розпізнавання райдужної оболонки ока, голосу. Серед усіх біометричних модальностей розпізнавання облич займає
особливе місце завдяки тому, що не потребує активної участі користувача і може працювати на відстані. \cite{Wang2021Survey}. Сучасні системи розпізнавання облич базуються на
глибоких згорткових нейронних мережах і демонструють точність понад 99\% в контрольованих умовах \cite{DHS2024}.
Архітектурно система розпізнавання облич складається з кількох ключових етапів: виявлення обличчя на зображенні або відеопотоці, нормалізація та вирівнювання
виявленої ділянки, екстракція ознак за допомогою нейронної мережі, та, нарешті, порівняння отриманого вектора ознак із базою даних еталонних зображень.
Результатом є рішення щодо підтвердження або відхилення ідентичності особи.
Системи розпізнавання облич активно впроваджуються в різноманітних сферах: контроль доступу до приміщень і пристроїв, верифікація особистості в банківській
сфері та при посадці на рейс, системи відеоспостереження та правоохоронна діяльність \cite{FacialRecognitionMarket2023}. За оцінками аналітиків, світовий ринок
цієї технології у 2023 році оцінювався у 6,3 млрд доларів США і прогнозується його зростання до 13,4 млрд доларів до 2028 року.
\subsection{Актуальність проблеми безпеки}
Незважаючи на високу точність у стандартних умовах, системи розпізнавання облич мають критичну вразливість — вони аналізують двовимірне зображення обличчя, а
не саму живу людину. Ця особливість відкриває можливість для так званих атак підміни (Presentation Attacks), коли зловмисник намагається обдурити систему,
пред'являючи їй фотографію, відео або тривимірну маску замість реального обличчя \cite{Yu2022Survey}.
Проблема набуває особливої гостроти з урахуванням широкого поширення систем розпізнавання облич у критично важливих застосуваннях. Несанкціонований доступ до мобільного пристрою,
банківського рахунку або охоронюваного об'єкта може мати серйозні фінансові та безпекові наслідки. При цьому для здійснення атаки підміни не потрібні спеціальні технічні знання —
достатньо якісної фотографії цільової особи, яку часто можна знайти у відкритих джерелах.
Дослідження у сфері протидії атакам на біометричні системи, відоме як Face Anti-Spoofing (FAS) або Presentation Attack Detection (PAD), є одним із найбільш активних напрямків сучасної
комп'ютерної біометрії. Кількість публікацій за ключовими словами «face anti-spoofing» та «face liveness detection» у наукометричних базах стрімко зростає з кожним роком, що свідчить про
актуальність та невирішеність проблеми \cite{Yu2022Survey}.
Метою даної роботи є дослідження та порівняльний аналіз сучасних методів на основі згорткових нейронних мереж для детекції 2D-атак підміни в системах біометричного розпізнавання облич, а
також практична реалізація та оцінка ефективності обраних архітектур на стандартному датасеті Replay-Attack \cite{Chingovska2012}.

% Список використаних джерел
\newpage
\begin{thebibliography}{99}

	\bibitem{Tome2015}
	Томе П., Марсель С. On the Vulnerability of Face Verification Systems to
	Hill-Climbing Attacks. \textit{IEEE TIFS}, 2015. Vol. 10, No. 5.

	\bibitem{Wang2021Survey}
	Ван М. та ін. Deep Face Recognition: A Survey.
	\textit{Neurocomputing}, 2021. Vol. 429.

	\bibitem{DHS2024}
	U.S. Department of Homeland Security. Biometric Technology Performance
	Report. Washington, DC: DHS, 2024.

	\bibitem{FacialRecognitionMarket2023}
	MarketsandMarkets. Facial Recognition Market — Global Forecast to 2028.
	2023.

	\bibitem{Yu2022Survey}
	Ю. Зітонг та ін. Deep Learning for Face Anti-Spoofing: A Survey.
	\textit{IEEE TPAMI}, 2023. Vol. 45, No. 5.

	\bibitem{Chingovska2012}
	Чінґовська І., Анжос А., Марсель С. On the Effectiveness of Local
	Binary Patterns in Face Anti-Spoofing. \textit{IEEE BIOSIG}, 2012.

\end{thebibliography}
\end{document}